\documentclass[11pt,oneside]{article}

\usepackage[T1]{fontenc}
\usepackage[utf8]{inputenc}
\usepackage{amsmath,amsthm,mathtools}
\usepackage{newtxtext,newtxmath}
\usepackage[final]{microtype}
\usepackage[a4paper,margin=30mm,headheight=14pt]{geometry}
\usepackage{booktabs,array,longtable,tabularx,multirow}
\usepackage{enumitem}
\usepackage{xcolor}
\usepackage[most]{tcolorbox}
\usepackage{fancyhdr}
\usepackage{hyperref}
\usepackage[nameinlink,capitalise,noabbrev]{cleveref}
\usepackage{listings}
\usepackage{url}
\usepackage{graphicx}
\usepackage{caption}
\usepackage{float}
\usepackage{siunitx}
\usepackage{pdflscape}

\definecolor{deepblue}{RGB}{31,57,94}
\definecolor{softblue}{RGB}{238,244,250}
\definecolor{softgray}{RGB}{246,246,246}
\definecolor{softgreen}{RGB}{239,248,241}
\definecolor{softamber}{RGB}{253,247,232}
\definecolor{softred}{RGB}{252,240,240}
\definecolor{rulegray}{RGB}{105,105,105}

\hypersetup{
  colorlinks=true,
  linkcolor=deepblue,
  citecolor=deepblue,
  urlcolor=deepblue,
  pdftitle={External-Prime Accumulation and the p-adic Newton Wall},
  pdfauthor={OpenAI GPT-5.6 Pro},
  pdfsubject={Progress report on the Alaoglu--Erdos integer-exponent conjecture}
}

\pagestyle{fancy}
\fancyhf{}
\fancyhead[L]{\small External-prime accumulation and the $p$-adic Newton wall}
\fancyhead[R]{\small 23 August 2026}
\fancyfoot[C]{\thepage}
\renewcommand{\headrulewidth}{0.4pt}

\setlength{\parindent}{0pt}
\setlength{\parskip}{0.55em}
\setlist{itemsep=0.25em,topsep=0.4em}
\allowdisplaybreaks

\newtheorem{theorem}{Theorem}[section]
\newtheorem{proposition}[theorem]{Proposition}
\newtheorem{lemma}[theorem]{Lemma}
\newtheorem{corollary}[theorem]{Corollary}
\newtheorem{conjecture}[theorem]{Conjecture}
\newtheorem{target}[theorem]{Target}
\theoremstyle{definition}
\newtheorem{definition}[theorem]{Definition}
\newtheorem{remark}[theorem]{Remark}
\newtheorem{example}[theorem]{Example}

\newtcolorbox{resultbox}[1]{
  enhanced,breakable,colback=softgreen,colframe=deepblue,
  boxrule=0.55pt,arc=1.5pt,left=6pt,right=6pt,top=5pt,bottom=5pt,
  title={#1},fonttitle=\bfseries
}
\newtcolorbox{statusbox}[1]{
  enhanced,breakable,colback=softblue,colframe=deepblue,
  boxrule=0.45pt,arc=1.5pt,left=6pt,right=6pt,top=5pt,bottom=5pt,
  title={#1},fonttitle=\bfseries
}
\newtcolorbox{warningbox}[1]{
  enhanced,breakable,colback=softamber,colframe=rulegray,
  boxrule=0.45pt,arc=1.5pt,left=6pt,right=6pt,top=5pt,bottom=5pt,
  title={#1},fonttitle=\bfseries
}
\newtcolorbox{deadbox}[1]{
  enhanced,breakable,colback=softred,colframe=rulegray,
  boxrule=0.45pt,arc=1.5pt,left=6pt,right=6pt,top=5pt,bottom=5pt,
  title={#1},fonttitle=\bfseries
}

\newcommand{\N}{\mathbb N}
\newcommand{\Nzero}{\mathbb N_0}
\newcommand{\Z}{\mathbb Z}
\newcommand{\Q}{\mathbb Q}
\newcommand{\R}{\mathbb R}
\newcommand{\C}{\mathbb C}
\newcommand{\Qp}{\mathbb Q_p}
\newcommand{\Cp}{\mathbb C_p}
\newcommand{\vp}{v_p}
\newcommand{\ord}{\operatorname{ord}}
\newcommand{\supp}{\operatorname{supp}}
\newcommand{\Sol}{\operatorname{Sol}}
\newcommand{\Span}{\operatorname{span}}
\newcommand{\Gauss}{\lVert\,\cdot\,\rVert}
\newcommand{\bigO}{\mathrm O}
\newcommand{\littleo}{\mathrm o}
\newcommand{\AEconj}{Alaoglu--Erd\H{o}s}

\lstdefinestyle{pseudocode}{
  basicstyle=\ttfamily\small,
  columns=fullflexible,
  frame=single,
  rulecolor=\color{rulegray},
  backgroundcolor=\color{softgray},
  breaklines=true,
  showstringspaces=false,
  xleftmargin=0.5em,
  xrightmargin=0.5em
}

\title{\textbf{External-Prime Accumulation and the $p$-Adic Newton Wall}\\[0.35em]
\large A nonduplicative progress report on the \AEconj{} integer-exponent conjecture}
\author{OpenAI GPT-5.6 Pro}
\date{23 August 2026}

\begin{document}
\maketitle

\begin{abstract}
Let
\[
 2^x\in\Z,\qquad 3^x\in\Z.
\]
The conjecture of Alaoglu and Erd\H{o}s asserts that $x\in\Z$.  The supplied unified report
already develops a very large corpus around the exact solution monoid, transcendence
reductions, determinant and interpolation methods, structural-prime residuals, $p$-adic
valuation lattices, Sturmian and carry encodings, canonical products, and geometric Newton
interpolation.  This report deliberately does not reproduce that corpus.  It opens a new
local direction at a prime $p\ge5$ dividing a hypothetical output.

The first result is a rigidity theorem for geometric sampling in a nonarchimedean field.  If
$|Q|_p<1$ and an analytic germ $F$ at zero satisfies
$F(Q^k)=R^k$ for every sufficiently large $k$, then $F(z)=z^d$ and $R=Q^d$ for some
$d\in\Nzero$.  A flexible two-row version says that
$F(\alpha Q^k)=\lambda F(Q^k)$ on an accumulating geometric ray forces the Taylor support of
$F$ into the exact eigenspectrum $\{n:\alpha^n=\lambda\}$.  In the \AEconj{} setting the edge
pairs are $(\alpha,\lambda)=(2,3)$ or $(3,2)$, so no nonzero analytic, meromorphic, or
finite-Puiseux germ can realize even one adjacent strip of the conditional integer grid.

The second result analyzes, at the same external prime, the geometric Newton coefficients
already identified in the companion report.  Their valuations are computed exactly.  The
leading coefficient of the unique degree-$N$ interpolation remainder forces every rational
polynomial interpolating the first $N+1$ native ray values to contain a coefficient with an
intrinsic denominator of order
\[
 p^{\frac12 v_p(Q)N^2+\mathrm O(N)}.
\]
On any off-ray unit multiple $uQ^k$, the Newton terms have valuations tending quadratically
to $-\infty$; eventually the last term uniquely dominates every partial sum.  Thus the
canonical Newton completion converges geometrically over $\C$, truncates exactly on its
native geometric ray in every completion, but diverges maximally over $\Qp$ at the adjacent
\AEconj{} row.  This is the \emph{$p$-adic Newton wall}.

The report then gives a Tate-algebra formulation of the edge operator
$F(z)\mapsto F(\alpha z)-\lambda F(z)$, an exact finite-strip defect lower bound, a normal-family
obstruction to uniformly overconvergent interpolants, a finite-rank $q$-difference research
interface, a computational audit, and a Lean formalization plan.  These are genuine new
lemmas and closure criteria, not a proof of the conjecture.  The remaining bridge is sharply
identified: one must derive nonarchimedean regularity, overconvergence, or finite-rank
$q$-difference structure from the global arithmetic data without assuming the desired power
branch is already regular at the accumulation point.
\end{abstract}

\tableofcontents
\newpage

\section{Mandate, status, and separation from the supplied report}

\subsection{The problem and the standard conditional notation}

The problem is
\begin{conjecture}[Alaoglu--Erd\H{o}s]\label{conj:AE}
For every real $x$,
\[
 2^x,3^x\in\Z \quad\Longrightarrow\quad x\in\Z.
\]
\end{conjecture}
It originates in the study of highly composite and related numbers
\cite{AE1944}.  It is a special instance of the still-open Four Exponentials boundary;
see, for example, Waldschmidt's survey \cite{Waldschmidt2023}.

Assume throughout the mathematical core that a nonintegral solution exists, and put
\begin{equation}\label{eq:MA}
 M=2^x\in\N,\qquad A=3^x\in\N,
 \qquad
 \theta=\frac{\log3}{\log2}.
\end{equation}
Then $A=M^\theta$ in the ordinary real sense.  The companion report
\cite{Unified2026} establishes much more conditional structure.  Only the small package
explicitly restated in \cref{sec:minimal-package} is imported here.

\subsection{What is new here}

The supplied report already contains the complex entire completion of the geometric Newton
interpolants and exact formulas for their coefficients.  What it does not contain is the
external-prime accumulation analysis carried out below.  The new material is:

\begin{enumerate}[label=\textbf{N\arabic*.}]
\item a geometric-ray classification theorem for analytic and meromorphic nonarchimedean
  germs;
\item an edge-law spectral theorem from zeros accumulating at the origin;
\item the corresponding exclusion of analytic, meromorphic, and finite-Puiseux realizations
  of an adjacent conditional \AEconj{} strip;
\item exact external-prime valuations of every geometric Newton coefficient, including the
  resonant valuation-tie correction;
\item a basis-independent quadratic lower bound for the $p$-adic Gauss norm of \emph{every}
  polynomial interpolant through a native geometric prefix;
\item exact off-ray valuations, an exact term-ratio formula, and eventual last-term dominance
  for Newton partial sums;
\item a Tate-algebra spectral-gap estimate for the dilation edge operator and a finite-strip
  quotient lower bound;
\item a compactness theorem showing that exact finite interpolants cannot have a uniformly
  overconvergent $p$-adic subsequence;
\item a concrete Robba-to-Tate and finite-rank $q$-difference research interface;
\item a reproducible exact-arithmetic audit and a staged Lean formalization plan.
\end{enumerate}

\begin{statusbox}{Status convention}
\textbf{Proved here} means that a complete conventional proof is included.  \textbf{Imported}
means that the statement is taken from the supplied unified report or a cited classical
source and used only as an input.  \textbf{Conditional target} means that the implication is
proved but the additional hypothesis has not been derived from a hypothetical
counterexample.  \textbf{Audit/no-go} records a rigorously identified obstruction and is not
presented as progress on the truth of the conjecture itself.
\end{statusbox}

\section{Public-state check as of 23 August 2026}
\label{sec:public-state}

The user's competitive premise is substantially grounded, but two headlines need precise
qualification.

\begin{enumerate}
\item Anthropic publicly launched Claude Fable 5 on 9 June 2026 and describes it as its most
  capable generally available model for long-horizon work \cite{AnthropicFable}.
\item The Jacobian development is an explicit Fable-assisted counterexample in dimension
  three, hence in every dimension at least three.  The two-dimensional Jacobian conjecture
  remains open.  Tao gives the explicit degree-seven map and a short verification
  \cite{TaoJacobian}.  Thus ``resolved the Jacobian conjecture'' is accurate only if it is
  understood as disproving the unrestricted conjecture, not settling the remaining planar
  case.
\item A curve with thirty exhibited independent rational points was posted to the ICARM
  leaderboard and attributed to Claude working with Levent Alp\H{o}ge and Ava Howell.  This
  proves an unconditional rank lower bound of $30$; an unconditional proof that the exact
  rank equals $30$ has not been claimed \cite{ICARM30,MathWorldRank30}.
\item The reported Hadamard construction supplies matrices for all twelve previously missing
  admissible orders below $2000$, including order $668$.  Epoch AI records the result and its
  human/AI attribution provisionally; OEIS records the twelve orders and attribution
  \cite{EpochHadamard,OEISHadamard}.  This closes that finite table, not the general Hadamard
  conjecture.
\end{enumerate}

Exact web searches for ``Alaoglu--Erd\H{o}s'', the integer-exponent statement, Fable, and Lean
found no public preprint, repository, announcement, or indexed social post describing the
rumored new route as of the date of this report.  I therefore treat the rival team's message
as unverified tactical information rather than as a mathematical input.  The search below is
independent of any claimed method.

\section{Minimal counterexample package and the external orientation}
\label{sec:minimal-package}

\subsection{Elementary consequences}

\begin{lemma}\label{lem:x-positive-trans}
Under \eqref{eq:MA}, a nonintegral solution $x$ is positive, irrational, and transcendental.
Moreover, $M$ is not a power of $2$, $A$ is not a power of $3$, and $2,M$ are
multiplicatively independent, as are $3,A$.
\end{lemma}

\begin{proof}
If $x<0$, then $0<2^x<1$, contradicting integrality; $x=0$ is integral.  If
$x=a/b\in\Q$ in lowest terms with $b>0$, then $2^a=M^b$.  Comparing the exponent of the
prime $2$ gives $b\mid a$, hence $b=1$.  Thus a nonintegral solution is irrational.  If it
were algebraic irrational, Gelfond--Schneider applied to $2^x=M\in\overline\Q$ would be a
contradiction, so it is transcendental.  The remaining statements follow immediately: for
example $M=2^m$ gives $x=m$, and $2^rM^s=1$ gives $r+sx=0$, forcing $r=s=0$ because $x$ is
irrational.
\end{proof}

The next fact is imported by the supplied corpus, but its short proof is included because it
is the control-sensitive hinge of this report.

\begin{lemma}[External-prime lemma]\label{lem:external-prime}
If \eqref{eq:MA} holds with $x\notin\Z$, then some prime $p\ge5$ divides $MA$.
\end{lemma}

\begin{proof}
Suppose instead that both $M$ and $A$ are $\{2,3\}$-smooth.  Write
\[
 M=2^a3^b,\qquad A=2^c3^d
 \qquad(a,b,c,d\in\Nzero).
\]
With $\theta=\log3/\log2$, the two logarithmic identities give
\[
 x=a+b\theta,\qquad x\theta=c+d\theta,
\]
so
\[
 b\theta^2+(a-d)\theta-c=0.
\]
The number $\theta$ is transcendental: it is not rational by unique factorization, and if it
were algebraic irrational then $2^\theta=3$ would contradict Gelfond--Schneider.  Therefore
$b=c=0$ and $a=d$, whence $x=a\in\Z$, a contradiction.
\end{proof}

\subsection{Oriented external-prime notation}

Choose an external prime $p\ge5$ from \cref{lem:external-prime}.  At least one of $M,A$ is
$p$-adically contracting.  We orient the data as follows:
\begin{equation}\label{eq:orientation}
 (Q,R;\alpha,\lambda,\gamma)=
 \begin{cases}
  (M,A;2,3,\theta),&p\mid M,\\
  (A,M;3,2,\theta^{-1}),&p\mid A.
 \end{cases}
\end{equation}
If $p$ divides both, either orientation may be used.  In both cases
\begin{equation}\label{eq:oriented-data}
 p\mid Q,
 \qquad Q^\gamma=R,
 \qquad \alpha^\gamma=\lambda,
 \qquad
 (Q^k)^\gamma=R^k,
 \qquad
 (\alpha Q^k)^\gamma=\lambda R^k.
\end{equation}
Set
\begin{equation}\label{eq:ma-local}
 m=\vp(Q)>0,
 \qquad
 a=\vp(R)\ge0.
\end{equation}
Because $p\ge5$, all of $2,3,\alpha,\lambda,$ and $\alpha-1$ are $p$-adic units.
The points $Q^k$ and $\alpha Q^k$ tend to $0$ in $\Qp$, giving a genuine
nonarchimedean accumulation point that is absent from the ordinary real geometry.

\begin{resultbox}{The control-sensitive feature}
For an integer control $x=d$, one has $(M,A)=(2^d,3^d)$, so no prime $p\ge5$ divides $MA$.
The oriented external-prime package \eqref{eq:orientation} therefore cannot even be formed for
controls.  This does not by itself prove the conjecture, but it passes the supplied report's
control-principle test: every theorem below visibly consumes a fact unavailable to the
integer controls.
\end{resultbox}

\section{Geometric rays at a nonarchimedean accumulation point}
\label{sec:analytic-rigidity}

\subsection{One-ray analytic classification}

The following elementary theorem is the cleanest new structural result in this report.

\begin{theorem}[Geometric-ray rigidity]\label{thm:ray-rigidity}
Let $K$ be a complete nonarchimedean field of characteristic zero.  Let $t,s\in K^\times$
with $0<|t|<1$, and let $F$ be analytic on a neighborhood of $0$.  If
\begin{equation}\label{eq:ray-values}
 F(t^k)=s^k
\end{equation}
for every sufficiently large integer $k$, then there is a unique $d\in\Nzero$ such that
\[
 s=t^d,
 \qquad
 F(z)=z^d
\]
on a neighborhood of $0$.
\end{theorem}

\begin{proof}
The values in \eqref{eq:ray-values} are nonzero, so $F$ is not identically zero.  Write
\[
 F(z)=z^dG(z),
 \qquad d=\ord_0F\in\Nzero,
 \qquad G(0)=c\ne0.
\]
Then
\begin{equation}\label{eq:u-powers-converge}
 \left(\frac{s}{t^d}\right)^k=G(t^k)\longrightarrow c.
\end{equation}
Put $u=s/t^d$.  Since the limit is nonzero, $|u|=1$.  A convergent sequence $u^k$ with
$|u|=1$ can converge to a nonzero limit only if $u=1$: indeed
\[
 u^{k+1}-u^k=u^k(u-1)\longrightarrow0,
\]
and $|u^k|=1$ forces $u-1=0$.  Hence $s=t^d$ and $G(t^k)=1$ for all sufficiently large
$k$.  The analytic function $G-1$ has infinitely many zeros tending to the interior point
$0$, so the nonarchimedean identity theorem gives $G\equiv1$.
\end{proof}

\begin{corollary}[One-ray obstruction for a hypothetical counterexample]
\label{cor:one-ray-AE}
In the oriented setting \eqref{eq:orientation}, no analytic germ $F$ at $0$ can satisfy
$F(Q^k)=R^k$ for every sufficiently large $k$.
\end{corollary}

\begin{proof}
By \cref{thm:ray-rigidity}, such a germ would give $R=Q^d$ for some $d\in\Nzero$.  In the
first orientation this is $3^x=(2^x)^d$, hence $3=2^d$ because $x>0$.  In the second it is
$2=3^d$.  Both are impossible.
\end{proof}

\begin{remark}[Newton-polygon reading]\label{rem:newton-slope}
The proof separates two obstructions.  First, an analytic germ has an integral order $d$ at
zero, so valuation comparison along the ray forces
\[
 \frac{\vp(R)}{\vp(Q)}=\frac am=d\in\Nzero.
\]
If $a/m$ is not an integer, the Newton slope already rules out analyticity.  If $a=dm$, the
unit quotient $R/Q^d$ must also have all of its powers converging; this forces the quotient to
be exactly $1$, not merely a root of unity or a principal unit.  Thus the valuation slope is
only the first layer; the unit mismatch is decisive in the resonant case.
\end{remark}

\subsection{The edge-law spectral theorem}

The two-row formulation no longer assumes that the native values themselves form a geometric
progression.

\begin{theorem}[Accumulating edge law]\label{thm:edge-law}
Let $K$ be as above, let $0<|t|<1$, and let $q,\mu\in K^\times$.  Suppose $F$ is analytic on
a disk $D$ centered at $0$ such that $qD\subseteq D$.  If
\begin{equation}\label{eq:edge-samples}
 F(qt^k)=\mu F(t^k)
\end{equation}
for infinitely many $k$ tending to infinity, then
\begin{equation}\label{eq:edge-spectrum}
 F(z)=\sum_{\substack{n\ge0\\q^n=\mu}}c_nz^n
\end{equation}
on $D$.  In particular, if $q^n\ne\mu$ for every $n\in\Nzero$, then $F\equiv0$.
\end{theorem}

\begin{proof}
The analytic function
\[
 H(z)=F(qz)-\mu F(z)
\]
has the zeros $t^k$ accumulating at the interior point $0$, so $H\equiv0$.  If
$F(z)=\sum_{n\ge0}c_nz^n$, coefficient comparison in $F(qz)=\mu F(z)$ gives
$(q^n-\mu)c_n=0$ for every $n$.
\end{proof}

\begin{corollary}[No analytic adjacent strip]\label{cor:no-analytic-strip}
In the oriented \AEconj{} setting, no nonzero analytic germ at $0$ can satisfy
\begin{equation}\label{eq:two-row-data}
 F(Q^k)=R^k,
 \qquad
 F(\alpha Q^k)=\lambda R^k
\end{equation}
for every sufficiently large $k$.
\end{corollary}

\begin{proof}
Apply \cref{thm:edge-law} with $(t,q,\mu)=(Q,\alpha,\lambda)$.  There is no
$n\in\Nzero$ with $2^n=3$ and no $n\in\Nzero$ with $3^n=2$, so $F\equiv0$, contradicting
$F(Q^k)=R^k\ne0$.
\end{proof}

\begin{remark}
The edge-law theorem is strictly more flexible than \cref{thm:ray-rigidity}: it uses only the
ratio between adjacent rows and does not use the formula $R^k$.  Conversely, the one-ray
theorem gives a contradiction before the adjacent row is introduced.  The two versions are
kept because they lead to different finite and functional-analytic targets later.
\end{remark}

\subsection{Continuity, meromorphic germs, and finite ramification}

Analyticity is genuinely stronger than continuity.  The support pattern of $Q$ and $R$
controls the first regularity wall.

\begin{proposition}[Continuity stratification]\label{prop:continuity}
Let $p\mid Q$.
\begin{enumerate}[label=(\roman*)]
\item If $p\nmid R$, no function continuous at $0$ can satisfy $F(Q^k)=R^k$ for all large
  $k$.
\item If $p\mid R$, the assignment $F(0)=0$, $F(Q^k)=R^k$ is continuous on the closed subset
  $\{0,Q^k:k\ge0\}$, but by \cref{cor:one-ray-AE} it has no analytic extension in the
  \AEconj{} setting.
\end{enumerate}
\end{proposition}

\begin{proof}
If $p\nmid R$ and $R^k$ converged in $\Qp$, then its limit would be nonzero and the same
successive-difference argument as in \cref{thm:ray-rigidity} would force $R=1$, impossible.
If $p\mid R$, then both $Q^k$ and $R^k$ tend to zero, so the displayed assignment is
continuous on that subset.
\end{proof}

This gives a useful local taxonomy.  An external prime lying in the symmetric difference of
$\supp(M)$ and $\supp(A)$ creates a \emph{continuity wall} in one orientation.  An external
prime common to both supports allows continuity but still creates an \emph{analyticity wall}.
The common-support case is therefore not locally invisible; it merely moves the obstruction
one rung up the regularity ladder.

\begin{theorem}[Meromorphic and finite-Puiseux rigidity]\label{thm:puiseux}
Let $K,t,s$ be as in \cref{thm:ray-rigidity}.
\begin{enumerate}[label=(\roman*)]
\item If a meromorphic germ $F$ at $0$ satisfies $F(t^k)=s^k$ for all sufficiently large
  $k$, then $F(z)=z^d$ and $s=t^d$ for some $d\in\Z$.
\item If, after a finite ramified cover $z=w^e$, a meromorphic germ realizes those values on
  a compatible branch, then $s^e=t^d$ for some $d\in\Z$.
\end{enumerate}
Consequently, the oriented \AEconj{} ray admits neither a meromorphic germ nor a finite-Puiseux
germ at $0$.
\end{theorem}

\begin{proof}
For (i), write the Laurent germ as $F(z)=z^dG(z)$ with $d\in\Z$ and $G(0)\ne0$; the proof of
\cref{thm:ray-rigidity} is unchanged.  For (ii), work in a complete algebraic closure and
choose $\tau$ with $\tau^e=t$.  Applying (i) to the pullback germ in $w$ gives
$s=\tau^d$, hence $s^e=t^d$.  In the \AEconj{} setting, $R^e=Q^d$ would imply
$3^e=2^d$ or $2^e=3^d$ after raising the defining real identities to the positive power
$1/x$, impossible by unique factorization.
\end{proof}

\begin{resultbox}{First frontier statement}
A hypothetical counterexample cannot be represented at its external accumulation point by
any ordinary analytic, meromorphic, or finitely ramified algebraic branch.  Therefore every
successful interpolation of the conditional grid must be essentially singular at zero,
infinitely ramified, nonlocal, or tied to a summation procedure that is not a convergent
$p$-adic germ.  The next sections show that the canonical geometric Newton interpolation
chooses precisely the last option: an archimedean sum whose $p$-adic terms explode.
\end{resultbox}

\section{The external-prime Newton wall}
\label{sec:newton-wall}

\subsection{The imported geometric Newton coefficients}

For integers $Q\ge2$ and $R\ge1$, let
\begin{equation}\label{eq:phi-cn}
 \phi_n(z)=\prod_{j=0}^{n-1}(z-Q^j),
 \qquad
 c_n(Q,R)=
 \frac{\prod_{j=0}^{n-1}(R-Q^j)}
 {Q^{\binom n2}\prod_{j=1}^{n}(Q^j-1)},
 \qquad c_0=1.
\end{equation}
The degree-$N$ Newton interpolant is
\begin{equation}\label{eq:PN}
 P_N^{Q,R}(z)=\sum_{n=0}^{N}c_n(Q,R)\phi_n(z).
\end{equation}
The supplied report proves the exact identity
\begin{equation}\label{eq:native-interpolation}
 P_N^{Q,R}(Q^k)=R^k
 \qquad(0\le k\le N),
\end{equation}
and constructs a locally uniformly convergent complex entire completion from the same
series.  Those facts are imported; the new point here is to inspect the rational coefficients
and summands in the completion $\Qp$ selected by an external prime.

The denominator in \eqref{eq:phi-cn} is obtained from
\[
 \prod_{j=0}^{n-1}(Q^n-Q^j)
 =Q^{\binom n2}\prod_{h=1}^n(Q^h-1).
\]
When $p\mid Q$, every $Q^h-1$ is a $p$-adic unit.  Thus the entire local problem is governed
by the numerator gaps $R-Q^j$ against the explicit quadratic denominator
$Q^{\binom n2}$.

\subsection{Exact valuation of every Newton coefficient}

We keep the oriented notation \eqref{eq:ma-local}: $m=\vp(Q)>0$ and
$a=\vp(R)\ge0$.
Because $(Q,R)$ is an oriented hypothetical \AEconj{} pair,
$R\ne Q^h$ for every $h\in\Nzero$; otherwise the defining positive real
identities would give $3=2^h$ or $2=3^h$.  Thus all numerator gaps below
are nonzero and their valuations are finite.

\begin{theorem}[Exact external-prime coefficient law]\label{thm:coeff-valuation}
For $n\ge1$, the following hold.

\begin{enumerate}[label=(\alph*)]
\item If $a=0$, then
\begin{equation}\label{eq:coeff-a0}
 \vp\bigl(c_n(Q,R)\bigr)
 =\vp(R-1)-m\binom n2.
\end{equation}

\item Suppose $a>0$.  Put
\begin{equation}\label{eq:qstar-epsilon}
 q_* =\left\lfloor\frac{a-1}{m}\right\rfloor,
 \qquad
 \varepsilon=
 \begin{cases}
  \vp\bigl(R-Q^{a/m}\bigr)-a,&m\mid a,\\
  0,&m\nmid a.
 \end{cases}
\end{equation}
Then, for every $n\ge q_*+2$,
\begin{equation}\label{eq:coeff-a-pos}
 \vp\bigl(c_n(Q,R)\bigr)
 =m\binom{q_*+1}{2}
  +a(n-1-q_*)+\varepsilon
  -m\binom n2.
\end{equation}
In particular,
\begin{equation}\label{eq:coeff-asymptotic}
 \vp\bigl(c_n(Q,R)\bigr)
 =-\frac m2n^2+\left(a+\frac m2\right)n+\bigO_{Q,R,p}(1).
\end{equation}
\end{enumerate}
\end{theorem}

\begin{proof}
Because $p\mid Q$, one has $\vp(Q^h-1)=0$ for $h\ge1$, and hence
\begin{equation}\label{eq:coeff-val-start}
 \vp(c_n)=\sum_{j=0}^{n-1}\vp(R-Q^j)-m\binom n2.
\end{equation}
If $a=0$, the term $j=0$ is $\vp(R-1)$, while for every $j\ge1$ the two terms $R$ and
$Q^j$ have distinct valuations $0$ and $mj>0$, so $\vp(R-Q^j)=0$.  This proves
\eqref{eq:coeff-a0}.

Now let $a>0$.  The $j=0$ gap $R-1$ is a unit.  For $j\ge1$ and $mj\ne a$, the
ultrametric inequality is strict and gives
\[
 \vp(R-Q^j)=\min(a,mj).
\]
The positive integers $j$ with $mj<a$ are exactly $1,\ldots,q_*$.  Their contribution is
$m\binom{q_*+1}{2}$.  Every later index contributes $a$, except that if $a=mh$ there is one
tie at $j=h$, whose excess above $a$ is exactly $\varepsilon$.  The range
$n\ge q_*+2$ includes the first index at or above the tie threshold.  Substitution into
\eqref{eq:coeff-val-start} gives \eqref{eq:coeff-a-pos}; expansion of the binomial coefficient
gives \eqref{eq:coeff-asymptotic}.
\end{proof}

\begin{remark}[The tie correction is finite but indispensable]
The term $\varepsilon$ records the case in which the valuation ray of $R$ meets the valuation
ray of $Q^j$.  It can be arbitrarily large for unrelated integer pairs.  In a hypothetical
\AEconj{} pair the equality $R=Q^{a/m}$ is impossible, so $\varepsilon$ is always finite, but it
cannot simply be discarded from exact finite statements.
\end{remark}

\begin{corollary}[Quadratic denominator escape]\label{cor:quadratic-denominator}
Let $D_n$ be the least positive integer such that $D_nc_n(Q,R)\in\Z_{(p)}$, where
$\Z_{(p)}$ is the localization away from $p$.  Then
\begin{equation}\label{eq:Dn-asymptotic}
 \vp(D_n)
 =\frac m2n^2-\left(a+\frac m2\right)n+\bigO_{Q,R,p}(1).
\end{equation}
In particular, the $p$-part of the denominator grows quadratically in $n$.
\end{corollary}

\subsection{The quadratic loss is intrinsic, not a Newton-basis artifact}

A change of interpolation basis cannot remove the local denominator.  The following theorem
applies to every polynomial that realizes the same finite data, regardless of its degree.

For a polynomial $P(z)=\sum b_jz^j\in\Qp[z]$, write
\[
 \lVert P\rVert_{\mathrm G}=\max_j|b_j|_p
\]
for the unit-disk Gauss norm.

\begin{theorem}[Universal Gauss-norm lower bound]\label{thm:universal-gauss}
Fix $N\ge0$.  Let $P\in\Qp[z]$ be any polynomial, of arbitrary degree, satisfying
\begin{equation}\label{eq:any-interpolant}
 P(Q^k)=R^k\qquad(0\le k\le N).
\end{equation}
Then
\begin{equation}\label{eq:gauss-lower}
 \lVert P\rVert_{\mathrm G}
 \ge |c_N(Q,R)|_p.
\end{equation}
Consequently, every rational polynomial satisfying \eqref{eq:any-interpolant} has at least
one coefficient whose reduced denominator has $p$-adic exponent at least
$\max\{0,-\vp(c_N)\}$.
\end{theorem}

\begin{proof}
Let
\[
 \Phi_{N+1}(z)=\prod_{k=0}^{N}(z-Q^k)\in\Z_p[z].
\]
It is monic.  Divide $P$ by $\Phi_{N+1}$ in $\Qp[z]$:
\[
 P=S\Phi_{N+1}+T,
 \qquad \deg T\le N.
\]
Because the divisor is monic with coefficients in $\Z_p$, the ordinary long-division
algorithm is nonexpanding for the Gauss norm: at each step one subtracts a multiple whose
coefficients have norm no larger than the current leading coefficient.  Hence
$\lVert T\rVert_{\mathrm G}\le\lVert P\rVert_{\mathrm G}$.

At every node $Q^k$, the divisor vanishes, so $T(Q^k)=P(Q^k)=R^k$.  By uniqueness of
interpolation in degree at most $N$, $T=P_N^{Q,R}$.  Its leading coefficient is precisely
$c_N(Q,R)$, because $\phi_N$ is monic and all preceding Newton terms have smaller degree.
Thus
\[
 |c_N|_p\le\lVert T\rVert_{\mathrm G}
 \le\lVert P\rVert_{\mathrm G}.
\]
\end{proof}

\begin{corollary}[No bounded exact polynomial family]\label{cor:no-bounded-polys}
For any sequence of rational polynomials $P_N$ satisfying
$P_N(Q^k)=R^k$ for $0\le k\le N$,
\[
 \log_p\lVert P_N\rVert_{\mathrm G}
 \ge \frac m2N^2-\bigO(N).
\]
Thus exact interpolation through longer native prefixes forces quadratic coefficient escape
in every polynomial model, not only in the canonical Newton one.
\end{corollary}

\begin{remark}
This lower bound is control-sensitive only through the choice of an external prime.  Similar
blow-up can occur at $2$ or $3$ for integer controls, but a prime $p\ge5$ contracting one of
the conditional outputs exists only in the nonintegral branch.  Any finishing argument still
has to use the fact that $p\nmid6$, not merely the existence of a contracting prime.
\end{remark}

\subsection{Exact valuation off the native ray}

Let $u\in\Z_p^\times$ and $k\in\Nzero$.  The point $uQ^k$ lies on the same valuation circle
as the native node $Q^k$ but is off the ray when $u\ne1$.

\begin{lemma}[Off-ray Newton-basis valuation]\label{lem:phi-offray}
For $n>k$,
\begin{equation}\label{eq:phi-offray}
 \vp\bigl(\phi_n(uQ^k)\bigr)
 =m\left(kn-\binom{k+1}{2}\right)+\vp(u-1).
\end{equation}
\end{lemma}

\begin{proof}
In
\[
 \phi_n(uQ^k)=\prod_{j=0}^{n-1}(uQ^k-Q^j),
\]
the factor $j=0$ is a unit when $k\ge1$; the formula also covers $k=0$ directly.  For
$1\le j<k$, the valuation is $mj$; at $j=k$ it is $mk+\vp(u-1)$; and for $j>k$ it is
$mk$.  Summing gives
\[
 m\sum_{j=1}^{k-1}j+mk+\vp(u-1)+mk(n-k-1)
 =m\left(kn-\frac{k(k+1)}2\right)+\vp(u-1).
\]
\end{proof}

Put
\begin{equation}\label{eq:Tn}
 T_n(z)=c_n(Q,R)\phi_n(z).
\end{equation}
Combining \cref{thm:coeff-valuation,lem:phi-offray} immediately gives:

\begin{corollary}[Quadratic off-ray explosion]\label{cor:term-explosion}
For fixed $u\in\Z_p^\times$, $u\ne1$, and fixed $k$, one has
\begin{equation}\label{eq:Tn-asymptotic}
 \vp\bigl(T_n(uQ^k)\bigr)
 =-\frac m2n^2+
 \left(a+m\left(k+\frac12\right)\right)n
 +\bigO_{Q,R,p,u,k}(1).
\end{equation}
Hence $T_n(uQ^k)$ does not tend to zero in $\Qp$; its $p$-adic absolute value tends to
infinity at a quadratic-exponential rate.
\end{corollary}

For the adjacent \AEconj{} row, $u=\alpha\in\{2,3\}$ and $p\ge5$, so
$\vp(u-1)=0$.  The formula is then especially clean.

\subsection{The exact term ratio and last-term dominance}

The two quadratic formulas above can be compressed into one exact first-order identity.

\begin{proposition}[Exact Newton-term ratio]\label{prop:term-ratio}
For every $n\ge0$ for which $T_n(z)\ne0$,
\begin{equation}\label{eq:term-ratio}
 \frac{T_{n+1}(z)}{T_n(z)}
 =\frac{(R-Q^n)(z-Q^n)}{Q^n(Q^{n+1}-1)}.
\end{equation}
For $z=uQ^k$, once $n>k$ and $mn>a$,
\begin{equation}\label{eq:term-ratio-vp}
 \vp\left(\frac{T_{n+1}(uQ^k)}{T_n(uQ^k)}\right)
 =a+mk-mn.
\end{equation}
For fixed complex $z$ in the ordinary absolute value,
\begin{equation}\label{eq:term-ratio-complex}
 \frac{T_{n+1}(z)}{T_n(z)}\longrightarrow\frac1Q.
\end{equation}
\end{proposition}

\begin{proof}
The quotient $c_{n+1}/c_n$ is
\[
 \frac{R-Q^n}{Q^n(Q^{n+1}-1)},
\]
and $\phi_{n+1}/\phi_n=z-Q^n$, giving \eqref{eq:term-ratio}.  At the external prime, for
$n>k$ one has $\vp(uQ^k-Q^n)=mk$, while $mn>a$ gives
$\vp(R-Q^n)=a$; $Q^{n+1}-1$ is a unit.  This proves
\eqref{eq:term-ratio-vp}.  In the ordinary absolute value, both numerator factors are
asymptotic to $-Q^n$, whereas the denominator is asymptotic to $Q^{2n+1}$.
\end{proof}

Equation \eqref{eq:term-ratio-complex} is the local reason for the complex entire convergence
proved in the companion report.  Equation \eqref{eq:term-ratio-vp} points in exactly the
opposite direction: its right side tends linearly to $-\infty$.

\begin{theorem}[Eventual last-term domination]\label{thm:last-term}
Fix $u\in\Z_p^\times$, $u\ne1$, and $k\in\Nzero$.  Then there exists $N_0$ such that for all
$N\ge N_0$,
\begin{equation}\label{eq:last-term}
 \vp\bigl(P_N^{Q,R}(uQ^k)\bigr)
 =\vp\bigl(T_N(uQ^k)\bigr).
\end{equation}
In particular,
\begin{equation}\label{eq:PN-offray-asymptotic}
 \vp\bigl(P_N^{Q,R}(uQ^k)\bigr)
 =-\frac m2N^2+
 \left(a+m\left(k+\frac12\right)\right)N+\bigO(1).
\end{equation}
\end{theorem}

\begin{proof}
By \eqref{eq:term-ratio-vp}, the valuations of successive terms are strictly decreasing once
$n>k+a/m$.  By \eqref{eq:Tn-asymptotic}, they tend to $-\infty$.  Therefore, for all
sufficiently large $N$, the final term $T_N$ has valuation strictly smaller than every
preceding term.  In a nonarchimedean field a sum with a unique term of least valuation has the
same valuation as that term, proving \eqref{eq:last-term}; the asymptotic follows from
\cref{cor:term-explosion}.
\end{proof}

\begin{corollary}[Adjacent-row repulsion]\label{cor:adjacent-repulsion}
In the oriented \AEconj{} setting, fix $k\in\Nzero$ and put $z=\alpha Q^k$.  Then, for all
sufficiently large $N$,
\begin{align}
 \vp\bigl(P_N^{Q,R}(z)-\lambda R^k\bigr)
 &=\vp\bigl(P_N^{Q,R}(z)\bigr)\label{eq:residual-equals-P}\\
 &=-\frac m2N^2+
 \left(a+m\left(k+\frac12\right)\right)N+\bigO(1).
 \label{eq:adjacent-residual-asymptotic}
\end{align}
Thus the native-ray interpolants become not merely inaccurate but quadratically
$p$-adically repelled from the prescribed adjacent-row integer.
\end{corollary}

\begin{proof}
The target $\lambda R^k$ has fixed finite valuation $ka$, because $\lambda$ is a $p$-adic
unit.  The valuation in \eqref{eq:PN-offray-asymptotic} tends to $-\infty$, so eventually it
is strictly smaller than $ka$.  The ultrametric inequality then gives
\eqref{eq:residual-equals-P}.
\end{proof}

\begin{resultbox}{The $p$-adic Newton wall}
The same rational Newton series exhibits three mutually compatible behaviors.
\begin{enumerate}[label=(\roman*)]
\item At a native node $z=Q^k$, every term with $n>k$ vanishes identically, so the value
  $R^k$ is a finite algebraic identity valid in every completion.
\item At a fixed complex point, the term ratio tends to $1/Q$, so the series converges
  geometrically and defines the entire function constructed in the supplied report.
\item At every off-ray point $z=uQ^k$ with $u$ a $p$-adic unit different from $1$, the term
  valuations tend to $-\infty$ quadratically, and the partial sums are eventually dominated
  by their last term.
\end{enumerate}
The native ray is therefore a set of exact truncation points surrounded, in the external
$p$-adic geometry, by maximal divergence.  This is not a numerical accident and not a
choice of summation order; it follows from the exact ratio \eqref{eq:term-ratio}.
\end{resultbox}

\subsection{Why the wall is not yet a contradiction}

The wall says that the canonical \emph{archimedean} completion of the native interpolation
data has no corresponding $p$-adic convergence off the ray.  It does not say that the
complex value at an off-ray rational point must equal a $p$-adic sum of the same terms.  An
infinite sum of rational numbers may converge in one completion and diverge in another; even
when it converges in both, the two limits need not be the same rational number.

At native nodes, equality is place-independent only because the series truncates.  At the
adjacent row it does not truncate, and \cref{cor:adjacent-repulsion} shows exactly how the
place dependence enters.  Any proof that transfers the complex entire completion into a
$p$-adic analytic germ without an additional arithmetic compatibility theorem would be
circular.

The product formula also does not close the gap.  A nonzero rational residual can have a
$p$-denominator of size $p^{\Theta(N^2)}$ while being only geometrically small in the
ordinary absolute value.  The elementary lower bound supplied by its denominator is then of
order $\exp(-\Theta(N^2))$, far below an archimedean error of order $\exp(-\Theta(N))$.
Quadratic denominator escape therefore defeats, rather than helps, a naive unit-locking
argument.

\begin{deadbox}{Closed shortcut: ``the entire function has rational data, so evaluate it $p$-adically''}
The Taylor coefficients of the complex completion are not obtained by a coefficientwise
finite rational operation; each fixed Taylor coefficient receives infinitely many Newton
contributions.  The complex convergence that defines those coefficients is itself
place-dependent.  Rational values at the native nodes come from exact truncation and do not
supply an adelic Taylor expansion.  This shortcut cannot be repaired without a separate
overconvergence or algebraicity theorem.
\end{deadbox}

\section{The Tate-algebra edge operator and finite-strip defects}
\label{sec:tate-edge}

The analytic contradiction in \cref{thm:edge-law} is qualitative.  On finite polynomials the
same mechanism is diagonal and quantitative.

\subsection{The edge spectrum}

For a radius $\rho>0$ and a polynomial $P(z)=\sum_{j=0}^db_jz^j\in\Qp[z]$, put
\begin{equation}\label{eq:weighted-gauss}
 \lVert P\rVert_\rho=\max_{0\le j\le d}|b_j|_p\rho^j.
\end{equation}
Define the edge operator
\begin{equation}\label{eq:edge-operator}
 \mathcal L_{\alpha,\lambda}P(z)=P(\alpha z)-\lambda P(z).
\end{equation}
Since $p\ge5$, $|\alpha|_p=|\lambda|_p=1$, and
\begin{equation}\label{eq:edge-diagonal}
 \mathcal L_{\alpha,\lambda}P
 =\sum_{j=0}^d(\alpha^j-\lambda)b_jz^j.
\end{equation}

\begin{definition}[Finite edge gap]\label{def:finite-gap}
For $d\in\Nzero$, set
\begin{equation}\label{eq:sigma-d}
 \sigma_{p,d}(\alpha,\lambda)
 =\min_{0\le j\le d}|\alpha^j-\lambda|_p.
\end{equation}
It is positive because no ordinary integer power of $2$ equals $3$, and no ordinary integer
power of $3$ equals $2$.
\end{definition}

\begin{proposition}[Finite spectral inequality]\label{prop:finite-spectral}
For every $P\in\Qp[z]$ of degree at most $d$,
\begin{equation}\label{eq:finite-spectral}
 \lVert\mathcal L_{\alpha,\lambda}P\rVert_\rho
 \ge \sigma_{p,d}(\alpha,\lambda)\lVert P\rVert_\rho.
\end{equation}
\end{proposition}

\begin{proof}
This is coefficientwise from \eqref{eq:edge-diagonal}:
\[
 \max_j |\alpha^j-\lambda|_p|b_j|_p\rho^j
 \ge\left(\min_j|\alpha^j-\lambda|_p\right)
      \max_j|b_j|_p\rho^j.
\]
\end{proof}

For the full Tate algebra, define
\begin{equation}\label{eq:sigma-infty}
 \sigma_p(\alpha,\lambda)
 =\inf_{j\ge0}|\alpha^j-\lambda|_p.
\end{equation}
The closure of the cyclic semigroup $\{\alpha^j:j\ge0\}$ is compact in $\Z_p^\times$, so
\begin{equation}\label{eq:closure-criterion}
 \sigma_p(\alpha,\lambda)>0
 \quad\Longleftrightarrow\quad
 \lambda\notin\overline{\{\alpha^j:j\ge0\}}.
\end{equation}
If $\lambda$ is not even a power of $\alpha$ modulo $p$, then every
$\alpha^j-\lambda$ is a $p$-adic unit and $\sigma_p=1$.

\begin{proposition}[Tate spectral gap]\label{prop:tate-gap}
If $\sigma_p(\alpha,\lambda)>0$, then on every Tate algebra on a disk stable under
$z\mapsto\alpha z$,
\begin{equation}\label{eq:tate-gap}
 \lVert\mathcal L_{\alpha,\lambda}F\rVert_\rho
 \ge\sigma_p(\alpha,\lambda)\lVert F\rVert_\rho.
\end{equation}
The operator is injective and has a bounded coefficientwise inverse on its image.
\end{proposition}

\begin{proof}
The proof of \cref{prop:finite-spectral} passes to convergent series.  Injectivity also follows
from the absence of an exact eigenvalue $\alpha^j=\lambda$.
\end{proof}

\begin{remark}[Resonant primes]
When $\lambda$ lies in the $p$-adic closure of powers of $\alpha$, the infinite gap vanishes,
but every finite gap \eqref{eq:sigma-d} remains positive.  The quantities
\[
 \mu_{p,d}(\alpha,\lambda)
 :=-\log_p\sigma_{p,d}(\alpha,\lambda)
 =\max_{0\le j\le d}\vp(\alpha^j-\lambda)
\]
measure the exact small-divisor loss.  This separates two local branches:
\emph{nonresonant} external primes, for which the edge operator is uniformly bounded below,
and \emph{resonant} external primes, where increasingly accurate congruences
$\alpha^j\equiv\lambda\pmod{p^r}$ may occur.  Any argument using a Tate inverse must keep this
branch distinction explicit.
\end{remark}

\subsection{A finite-strip quotient lower bound}

Let $K\le N$ and define the tail-node polynomial
\begin{equation}\label{eq:PhiKN}
 \Phi_{K,N}(z)=\prod_{k=K}^{N}(z-Q^k).
\end{equation}
If a polynomial $P$ satisfies the adjacent-row identity at those nodes, then
$\mathcal L_{\alpha,\lambda}P$ is divisible by $\Phi_{K,N}$.

\begin{theorem}[Finite-strip defect theorem]\label{thm:finite-strip}
Let $P\in\Qp[z]$ have degree at most $d$ and suppose
\begin{equation}\label{eq:finite-strip-hyp}
 P(\alpha Q^k)=\lambda P(Q^k)
 \qquad(K\le k\le N).
\end{equation}
Write
\begin{equation}\label{eq:defect-factorization}
 \mathcal L_{\alpha,\lambda}P(z)
 =\Phi_{K,N}(z)\,S(z).
\end{equation}
Fix $0<\rho<1$ such that $|Q|_p^K\le\rho$.  If additionally
$P(Q^K)=R^K$, then
\begin{equation}\label{eq:defect-lower}
 \lVert S\rVert_\rho
 \ge
 \sigma_{p,d}(\alpha,\lambda)
 |R|_p^K\rho^{-(N-K+1)}.
\end{equation}
\end{theorem}

\begin{proof}
The roots in \eqref{eq:finite-strip-hyp} give the factorization
\eqref{eq:defect-factorization}.  The weighted Gauss norm is multiplicative, and because
$|Q|_p^k\le|Q|_p^K\le\rho$ for every $k\ge K$,
\begin{equation}\label{eq:Phi-norm}
 \lVert\Phi_{K,N}\rVert_\rho
 =\prod_{k=K}^{N}\max(\rho,|Q|_p^k)
 =\rho^{N-K+1}.
\end{equation}
Evaluation inside the disk gives
\[
 |R|_p^K=|P(Q^K)|_p\le\lVert P\rVert_\rho.
\]
Apply \cref{prop:finite-spectral} and multiplicativity:
\[
 \rho^{N-K+1}\lVert S\rVert_\rho
 =\lVert\mathcal L_{\alpha,\lambda}P\rVert_\rho
 \ge\sigma_{p,d}\lVert P\rVert_\rho
 \ge\sigma_{p,d}|R|_p^K.
\]
\end{proof}

\begin{corollary}[No tame defect quotient]\label{cor:no-tame-defect}
Let $P_N$ be the unique polynomial of degree at most $2N+1$ interpolating both rows
\[
 P_N(Q^k)=R^k,
 \qquad
 P_N(\alpha Q^k)=\lambda R^k
 \qquad(0\le k\le N).
\]
Fix $K$ with $|Q|_p^K\le\rho<1$.  In the factorization
\[
 \mathcal L_{\alpha,\lambda}P_N=\Phi_{K,N}S_{K,N},
\]
one has
\begin{equation}\label{eq:SKN-growth}
 \lVert S_{K,N}\rVert_\rho
 \ge
 \sigma_{p,2N+1}(\alpha,\lambda)
 |R|_p^K\rho^{-(N-K+1)}.
\end{equation}
At a nonresonant external prime, the quotient norm therefore grows at least geometrically
like $\rho^{-N}$.
\end{corollary}

The zeros forced by a long adjacent strip make the defect small on a strict disk, but the
edge spectral inequality prevents the whole polynomial from becoming small.  The quotient
$S_{K,N}$ must absorb the difference.  This is the finite analogue of the essential
singularity forced by \cref{cor:no-analytic-strip}.

\begin{target}[Subexponential defect quotient]\label{target:defect-quotient}
Find an arithmetic construction of exact finite two-row interpolants for which, at some
external prime and some fixed $\rho<1$,
\[
 \log\lVert S_{K,N}\rVert_\rho=o(N)
 \quad\text{and}\quad
 -\log\sigma_{p,\deg P_N}(\alpha,\lambda)=o(N).
\]
Then \eqref{eq:defect-lower} is impossible for large $N$, proving the conjecture in that
branch.
\end{target}

The first requirement is the substantive one.  Canonical polynomial interpolation generally
has much larger local height, and \cref{thm:universal-gauss} shows that no native-ray exact
interpolant can be uniformly bounded.  Nevertheless, \cref{target:defect-quotient} is a
precise finite target: it asks for control of the \emph{edge quotient}, not of every
coefficient of the interpolant.

\section{Overconvergence, compactness, and a Robba-to-Tate bridge}
\label{sec:overconvergence}

\subsection{A normal-family obstruction}

The local classification theorem can be reached from finite data by compactness, provided
one has genuine overconvergence.
For a convergent series $F(z)=\sum_{j\ge0}b_jz^j$ on the closed disk of radius
$\rho$, write $\lVert F\rVert_\rho=\sup_j |b_j|_p\rho^j$.

\begin{theorem}[No uniformly overconvergent interpolation family]
\label{thm:no-overconvergent-family}
Let $\rho>1$.  Suppose $F_N(z)=\sum_{j\ge0}b_{N,j}z^j$ is analytic on the closed
$p$-adic disk of radius $\rho$, and suppose
\begin{equation}\label{eq:FN-prefix}
 F_N(Q^k)=R^k\qquad(0\le k\le N).
\end{equation}
Then the norms $\lVert F_N\rVert_\rho$ are unbounded as $N\to\infty$.
The same conclusion holds if \eqref{eq:FN-prefix} is replaced by the two-row prefix.
\end{theorem}

\begin{proof}
Assume $\lVert F_N\rVert_\rho\le C$.  Then
$|b_{N,j}|_p\le C\rho^{-j}$.  Closed bounded subsets of $\Qp$ are compact, so a diagonal
subsequence makes every coefficient $b_{N,j}$ converge.  The geometric bound on the tails
makes that subsequence converge uniformly on the unit disk to an analytic function
$F(z)=\sum b_jz^j$.  For each fixed $k$, the equality $F_N(Q^k)=R^k$ holds eventually, so
uniform convergence gives $F(Q^k)=R^k$.  This contradicts
\cref{cor:one-ray-AE}.  The two-row version is immediate as well.
\end{proof}

\begin{remark}
Boundedness on the unit disk alone is not enough for this compactness argument: coefficient
bounds need not give a uniformly small tail.  The strict radius gain $\rho>1$ is exactly the
overconvergence input.  This is why a Robba-ring or punctured-disk construction is not yet a
solution; it must be shown to descend through the puncture with a quantitative radius margin.
\end{remark}

\begin{corollary}[Coefficient escape for any overconvergent model]
Any sequence of exact finite-prefix models analytic on one common disk of radius
$\rho>1$ must have
\[
 \limsup_{N\to\infty}\lVert F_N\rVert_\rho=\infty.
\]
For polynomial models, \cref{thm:universal-gauss} gives the stronger explicit lower bound
$p^{mN^2/2-\bigO(N)}$ already on the unit disk.
\end{corollary}

\subsection{The removable-singularity interface}

The grid values themselves naturally live on a punctured set converging to zero.  It is easy
to imagine functions on a punctured disk that realize arbitrary values there while having an
essential singularity at zero.  The useful theorem would not be another interpolation
existence statement, but a descent theorem that makes the singularity removable.

\begin{target}[Robba-to-Tate descent]\label{target:robba-tate}
Construct, from the integer grid alone, a function or finite-rank difference module over a
Robba ring at an external prime with all three properties:
\begin{enumerate}[label=(\roman*)]
\item it realizes the native and adjacent row values in \eqref{eq:two-row-data};
\item its slopes, denominators, or Frobenius structure satisfy an arithmetic boundedness
  condition derived from $M,A\in\Z$;
\item that boundedness forces descent to a Tate algebra, a meromorphic germ, or a finite
  Puiseux extension at zero.
\end{enumerate}
Any such descent proves the conjecture by
\cref{cor:no-analytic-strip,thm:puiseux}.
\end{target}

This target is deliberately modular.  It does not demand that the power branch itself be
postulated $p$-adically.  It asks for an arithmetic object first and regularity second.  The
Newton wall says what the descent theorem must defeat: the obvious one-axis completion has
quadratically escaping denominators and no off-ray $p$-adic summation.

\subsection{A coefficient-growth criterion}

Let $F(z)=\sum_{n\ge n_0}b_nz^n$ be a formal Laurent series over $\Qp$.  It defines a
meromorphic germ on a punctured neighborhood of zero whenever the positive tail has positive
radius of convergence and the negative tail is finite.  By \cref{thm:puiseux}, no such series
can realize the oriented ray.

Thus any formal object carrying the data must satisfy at least one of:
\begin{enumerate}
\item infinitely many negative powers;
\item zero $p$-adic radius for the positive tail;
\item exponents with unbounded denominators, so that no finite ramified cover turns the object
  into a Laurent series;
\item failure of the formal values to agree with the integer samples under convergent
  evaluation.
\end{enumerate}
This four-way escape list is useful when auditing proposed constructions.  For example, a
$G$-function-style series with algebraic coefficients, controlled denominators, and positive
$p$-adic radius at almost every prime would immediately be incompatible with an external
\AEconj{} ray, provided its sample equalities were algebraic identities rather than purely
archimedean limits.

\begin{target}[Adelic positive-radius criterion]\label{target:adelic-radius}
Produce a series $F(z)=\sum b_nz^n$ with coefficients in one number field such that:
\begin{enumerate}[label=(\roman*)]
\item the equalities $F(Q^k)=R^k$ and $F(\alpha Q^k)=\lambda R^k$ are certified by finite
  algebraic identities or by a summation theorem valid in every completion;
\item at every prime outside a fixed set, and in particular at an external $p$, the
  coefficients have at most exponential denominator growth, so the $p$-adic radius is
  positive.
\end{enumerate}
Then the external-prime edge theorem closes the conjecture.
\end{target}

The first clause is essential.  Complex equality of an infinite sum to an integer is not an
adelic identity.

\section{A finite-rank \texorpdfstring{$q$}{q}-difference route}
\label{sec:q-difference}

\subsection{Why two independent dilations are structurally promising}

A genuine interpolation of the conditional grid would carry two dilation laws,
\begin{equation}\label{eq:two-dilations}
 F(2z)=3F(z),
 \qquad
 F(Mz)=AF(z).
\end{equation}
Under a nonintegral solution, $2$ and $M$ are multiplicatively independent by
\cref{lem:x-positive-trans}.  Systems of linear $q$-difference equations for two
multiplicatively independent dilation parameters are highly rigid under rationality,
finite-rank, and consistency hypotheses.  Results of B\'ezivin and Boutabaa, and the later
unified treatment of Sch\"afke and Singer, show that common formal or meromorphic solutions
of sufficiently regular independent difference systems often reduce to rational or constant
normal forms \cite{Bezivin1992,BezivinBoutabaa1992,SchaefkeSinger2019}.

The present scalar equations are already consistent: the constants commute and
$3A=A3$.  Their formal Laurent obstruction is immediate.  If
$F(z)=\sum_{n\in\Z}c_nz^n$ has only finitely many negative powers, then the first equation in
\eqref{eq:two-dilations} gives
\[
 (2^n-3)c_n=0
\]
for every $n$, hence $F=0$.  On the universal cover of $\C^\times$, however, the nonzero
solution $z^\theta$ exists.  The gap is exactly local regularity and monodromy, not
consistency.

\subsection{Why existing two-\texorpdfstring{$q$}{q} theorems do not apply automatically}

The integer table gives values at a discrete semigroup, not a finite-dimensional difference
module over $\overline\Q(z)$.  The complex geometric Newton completion satisfies a one-axis
\emph{inhomogeneous} equation with a $q$-product defect; it does not satisfy both homogeneous
laws in \eqref{eq:two-dilations}.  The external-prime Newton wall further shows that this
completion has no off-ray $p$-adic germ from which one could read a local difference module.

Therefore it would be an overstatement to invoke a two-$q$ rationality theorem directly.
One first needs a finite-rank closure of the arithmetic defects.

\begin{target}[Finite-rank edge cocycle]\label{target:finite-rank-cocycle}
Construct a vector $Y(z)$ of finitely many arithmetic interpolation functions and rational
matrices $B_2(z),B_M(z)\in\operatorname{GL}_r(\overline\Q(z))$ such that
\begin{equation}\label{eq:finite-rank-system}
 Y(2z)=B_2(z)Y(z),
 \qquad
 Y(Mz)=B_M(z)Y(z),
\end{equation}
with the consistency identity
\begin{equation}\label{eq:consistency}
 B_2(Mz)B_M(z)=B_M(2z)B_2(z),
\end{equation}
and with one coordinate of $Y$ carrying the exact conditional integer grid.  Prove that the
system is regular singular or overconvergent at an external prime.  A reduction theorem for
consistent independent $q$-difference systems would then place the coordinate in a
Laurent/Puiseux class, contradicted by \cref{thm:puiseux}.
\end{target}

This target is narrower than ``find an analytic interpolation.''  Finite rank can sometimes
be proved from recurrences among defects even when no scalar completion is regular.  The
companion report's exact endpoint defect and two-axis connection formulas provide raw
material, but their current defect family contains nonrational infinite $q$-products and is
not known to close in finite rank.

\subsection{A practical search protocol}

A computational search for \cref{target:finite-rank-cocycle} should not fit values blindly.
It should:
\begin{enumerate}
\item choose a small list of canonical one-axis completions and endpoint defects from the
  supplied report;
\item apply both dilation operators repeatedly and reduce every expression in a fixed
  symbolic basis of $q$-products and Newton kernels;
\item compute the rank of the span over $\Q(z)$ at increasing truncation order;
\item reject bases whose rank grows with the truncation order;
\item for any apparently stable rank, solve the exact rational-function consistency equations
  and verify them symbolically;
\item only then analyze local slopes at primes $p\mid MA$, with external primes separated
  from $2$ and $3$.
\end{enumerate}
A stable finite rank would be a substantive new structure theorem.  Failure with rank growth
would close another route without pretending that numerical near-dependence is a proof.

\section{Adelic and control audits}
\label{sec:audits}

\subsection{The integer controls remain an important trap}

For a control exponent $x=d\in\N$, the real branch $z^\theta$ still has no $2$-adic analytic
germ interpolating $(2^{dk},3^{dk})$.  Therefore the mere statement ``the irrational power
branch is not $p$-adically analytic'' is control-blind.  The external-prime program becomes
potentially discriminating only when a construction uses $p\nmid6$ to obtain regularity,
bounded denominators, an edge spectral gap, or a finite-rank local module.  Any proposed
bridge should answer the concrete question:

\begin{quote}
Which step works at a prime $p\ge5$ dividing a hypothetical output but has no analogue at the
structural primes $2$ and $3$ of the integer controls?
\end{quote}

The exact valuation formulas provide one possible answer: at an external prime, the
quadratic denominator lies outside $\Z[1/6]$ and directly interacts with an edge operator
whose two multipliers are units.  At present that interaction diagnoses singularity; it does
not yet reverse it into regularity.

\subsection{No contradiction from the product formula at current scales}

The canonical off-ray term has
\[
 \log|T_n(uQ^k)|_p
 =\frac m2n^2\log p+\bigO(n),
\]
while in the ordinary absolute value its ratio tends to $1/Q$, so
$\log|T_n|_\infty=-n\log Q+\bigO(1)$ at fixed $z$ away from accidental zeros.  These two
scales do not oppose one another: the remaining finite places absorb the quadratic balance.
A useful product-formula object would need either
\begin{enumerate}
\item quadratic archimedean smallness matching the external denominator;
\item cancellation of the external denominator in the \emph{reduced} rational residual; or
\item simultaneous positive valuation at enough other places to exceed the global height.
\end{enumerate}
None follows from the one-axis Newton formula.

\subsection{Sampling is not regularity}

The points $Q^k$ accumulate at zero $p$-adically, but a function specified only on those
points need not be continuous, let alone analytic.  Integrality of the values does not change
that.  \Cref{thm:ray-rigidity,thm:edge-law} should therefore be read as
\emph{closure theorems}: once a regular germ exists, the conjecture closes immediately.  The
construction of that germ, or of a finite-rank object forced to become one, is the missing
arithmetic theorem.

\begin{warningbox}{Logical boundary}
No argument in this report infers $p$-adic continuity from real continuity, infers a $p$-adic
sum from a complex sum, or treats integer-valued sampling as an analytic extension theorem.
Those three inferences are precisely where several attractive one-page ``proofs'' fail.
\end{warningbox}

\section{Exact computational audit}
\label{sec:computational-audit}

The new valuation statements are symbolic, but an exact audit is useful for detecting index
errors, especially in the valuation-tie case.  The audit uses Python's
\texttt{fractions.Fraction}; no floating-point arithmetic is involved.  For each test case it
constructs $c_n$, $\phi_n(uQ^k)$, the Newton term $T_n$, and the partial sum $P_N$ as exact
rational numbers, then compares directly computed valuations with the formulas in
\cref{thm:coeff-valuation,lem:phi-offray}.

The test script used for this report has SHA-256 digest
\begin{center}
\ttfamily\small
20458cc7e3174c8a7d30f009e8259eba1f6629b044fdc33a997b495a6bdd49fe.
\end{center}
Its complete source is reproduced in \cref{app:audit-code}.

\subsection{Case I: unit output at the contracting prime}

Take
\[
 Q=35,\qquad R=22,\qquad p=5,\qquad u=2,\qquad k=2.
\]
Here $m=1$, $a=0$, and $\vp(R-1)=0$.  Thus
\[
 \vp(c_n)=-\binom n2,
 \qquad
 \vp\bigl(\phi_n(2Q^2)\bigr)=2n-3\quad(n>2),
\]
so
\begin{equation}\label{eq:case1-term}
 \vp(T_n(2Q^2))=-\frac{(n-2)(n-3)}2
 \qquad(n>2).
\end{equation}
The exact audit is:

\begin{table}[H]
\centering
\caption{Case I: exact valuations.}
\label{tab:case1}
\begin{tabular}{r@{\qquad}r@{\qquad}r@{\qquad}r@{\qquad}r}
\toprule
$n$ & $\vp(c_n)$ & $\vp(\phi_n)$ & $\vp(T_n)$ & $\vp(P_n)$\\
\midrule
0&0&0&0&0\\
1&0&0&0&0\\
2&-1&1&0&0\\
3&-3&3&0&1\\
4&-6&5&-1&-1\\
5&-10&7&-3&-3\\
6&-15&9&-6&-6\\
7&-21&11&-10&-10\\
8&-28&13&-15&-15\\
9&-36&15&-21&-21\\
10&-45&17&-28&-28\\
11&-55&19&-36&-36\\
12&-66&21&-45&-45\\
\bottomrule
\end{tabular}
\end{table}

From $n=4$ onward, the final term is already the unique term of least valuation, so the
partial sum has exactly the predicted last-term valuation.

\subsection{Case II: a genuine valuation tie}

Take
\[
 Q=10,\qquad R=110,\qquad p=5,\qquad u=2,\qquad k=1.
\]
Now $m=a=1$, and the tie occurs at $j=1$:
\[
 \varepsilon=\vp(110-10)-1=\vp(100)-1=1.
\]
Formula \eqref{eq:coeff-a-pos} gives, for $n\ge2$,
\[
 \vp(c_n)=n-\binom n2.
\]
The audit confirms that omitting $\varepsilon$ would shift every coefficient valuation from
$n=2$ onward.

\begin{table}[H]
\centering
\caption{Case II: exact valuations with a nonzero tie correction.}
\label{tab:case2}
\begin{tabular}{r@{\qquad}r@{\qquad}r@{\qquad}r@{\qquad}r}
\toprule
$n$ & $\vp(c_n)$ & $\vp(\phi_n)$ & $\vp(T_n)$ & $\vp(P_n)$\\
\midrule
0&0&0&0&0\\
1&0&0&0&1\\
2&1&1&2&1\\
3&0&2&2&1\\
4&-2&3&1&1\\
5&-5&4&-1&-1\\
6&-9&5&-4&-4\\
7&-14&6&-8&-8\\
8&-20&7&-13&-13\\
9&-27&8&-19&-19\\
10&-35&9&-26&-26\\
11&-44&10&-34&-34\\
12&-54&11&-43&-43\\
\bottomrule
\end{tabular}
\end{table}

\subsection{Case III: integral valuation slope but unit mismatch}

Take
\[
 Q=25,\qquad R=75,\qquad p=5,\qquad u=2,\qquad k=1.
\]
Here $m=a=2$, so $a/m=1$ is an integral candidate order.  Nevertheless
$R/Q=3\ne1$, exactly the unit obstruction in \cref{rem:newton-slope}.  The tie excess is zero,
and
\[
 \vp(c_n)=-(n-1)(n-2)\qquad(n\ge2).
\]

\begin{table}[H]
\centering
\caption{Case III: an integral slope does not prevent quadratic escape.}
\label{tab:case3}
\begin{tabular}{r@{\qquad}r@{\qquad}r@{\qquad}r@{\qquad}r}
\toprule
$n$ & $\vp(c_n)$ & $\vp(\phi_n)$ & $\vp(T_n)$ & $\vp(P_n)$\\
\midrule
0&0&0&0&0\\
1&0&0&0&2\\
2&0&2&2&2\\
3&-2&4&2&2\\
4&-6&6&0&0\\
5&-12&8&-4&-4\\
6&-20&10&-10&-10\\
7&-30&12&-18&-18\\
8&-42&14&-28&-28\\
9&-56&16&-40&-40\\
10&-72&18&-54&-54\\
11&-90&20&-70&-70\\
12&-110&22&-88&-88\\
\bottomrule
\end{tabular}
\end{table}

All formula checks through $n=12$ pass exactly in the three branches.  The cases are not
intended as candidate \AEconj{} outputs; they are local unit tests covering $a=0$, a positive
valuation with nonzero tie excess, and an integral slope with unit mismatch.

\subsection{What to compute next}

The next computational experiment should not scan arbitrary integer pairs for a false
counterexample.  It should test the finite-rank and defect-quotient targets:

\begin{enumerate}
\item For many pairs $(Q,R)$ and external primes $p\mid Q$, construct the exact two-row
  interpolant through $Q^k,\alpha Q^k$ for $0\le k\le N$.
\item Factor $\mathcal L_{\alpha,\lambda}P_N$ by the largest forced tail product
  $\Phi_{K,N}$ and record the weighted Gauss norm of the quotient.
\item Separate primes by the resonance statistic $\mu_{p,d}(\alpha,\lambda)$.
\item Compare actual pairs arising from the supplied finite search with integer controls and
  with random pairs having the same valuation profile.
\item Apply repeated $2$- and $Q$-dilations to the canonical defect family and compute exact
  ranks over $\Q(z)$, looking for stabilization rather than numerical singular values.
\end{enumerate}

A stable finite rank or subexponential defect quotient would be a positive discovery.  Growth
linear in $N$ for the rank, or geometric growth saturating \eqref{eq:SKN-growth}, would be a
clean negative closure.

\section{Lean formalization plan}
\label{sec:lean-plan}

The new package splits naturally into finite algebra, elementary valuation combinatorics, and
nonarchimedean analysis.  The first two layers are substantially easier to formalize than the
full analytic identity theorem and should be completed first.

\subsection{Proposed module graph}

\begingroup\small
\begin{longtable}{@{}p{0.35\textwidth}p{0.16\textwidth}p{0.41\textwidth}@{}}
\caption{Proposed Lean modules and trust boundary.}\label{tab:lean-modules}\\
\toprule
Module & Initial status & Main content\\
\midrule
\endfirsthead
\toprule
Module & Initial status & Main content\\
\midrule
\endhead
\texttt{ExternalPrimeOrientation} & finite / elementary & The external-prime lemma; orientation data; unit facts for $p\ge5$; impossibility of $R=Q^d$ in either \AEconj{} orientation.\\
\texttt{QNewtonCoeffValuation} & Lean-ready & Formula \eqref{eq:coeff-val-start}; $a=0$ branch; threshold $q_*$; unique tie and excess $\varepsilon$; exact formula \eqref{eq:coeff-a-pos}.\\
\texttt{OffRayNewtonBasis} & Lean-ready & Factor-by-factor proof of \eqref{eq:phi-offray}; exact term ratio \eqref{eq:term-ratio}; eventual negative ratio valuation.\\
\texttt{NewtonLastTermDominance} & Lean-ready & Strictly decreasing valuations, unique-minimum valuation of a finite sum, and \eqref{eq:PN-offray-asymptotic}.\\
\texttt{InterpolationGaussLowerBound} & finite polynomial & Monic division over $\Qp$; nonexpansion of Gauss norm; uniqueness of the interpolation remainder; \cref{thm:universal-gauss}.\\
\texttt{EdgeOperatorSpectrum} & finite polynomial & Diagonal coefficient action, finite gap $\sigma_{p,d}$, and inequality \eqref{eq:finite-spectral}.\\
\texttt{FiniteStripDefect} & finite polynomial & Root factorization, weighted Gauss norm of $\Phi_{K,N}$, and \cref{thm:finite-strip}.\\
\texttt{GeometricRayAnalytic} & analysis layer & Order of an analytic germ, convergence of powers, the nonarchimedean identity theorem, and \cref{thm:ray-rigidity}.\\
\texttt{EdgeLawAnalytic} & analysis layer & Strassmann/identity theorem for accumulating zeros and coefficient-spectrum theorem \cref{thm:edge-law}.\\
\texttt{OverconvergentCompactness} & analysis/topology & Diagonal compactness of bounded coefficient balls, uniform-tail convergence, and \cref{thm:no-overconvergent-family}.\\
\bottomrule
\end{longtable}
\endgroup

\subsection{Finite theorem signatures}

The valuation layer can be stated without committing to a particular analytic API.  A
possible pseudocode interface is:

\begin{lstlisting}[style=pseudocode,language={}] 
-- Q, R, p, n are natural numbers; p is prime and p | Q.
def qNewtonCoeff (Q R n : N) : Q :=
  (prod j in range n, (R : Z) - Q^j) /
  (Q^(n.choose 2) * prod j in Icc 1 n, (Q^j : Z) - 1)

theorem padicVal_qNewtonCoeff_of_val_eq_zero
    (hpQ : p | Q) (hR : padicValNat p R = 0) (hn : 1 <= n) :
  padicValRat p (qNewtonCoeff Q R n) =
    padicValInt p (R - 1) - padicValNat p Q * n.choose 2

theorem padicVal_qNewtonCoeff_eventually
    (hpQ : p | Q) (ha : 0 < padicValNat p R)
    (hn : qStar Q R p + 2 <= n) :
  padicValRat p (qNewtonCoeff Q R n) =
    exactLocalFormula Q R p n

theorem padicVal_qNewtonTerm_offRay
    (hpQ : p | Q) (hpu : not p | u) (hk : k < n) :
  padicValRat p
    (qNewtonCoeff Q R n * prod j in range n, (u*Q^k : Z) - Q^j) =
  exactTermFormula Q R p u k n
\end{lstlisting}

The actual Mathlib valuation types may favor additive valuations into
$\mathbb Z_{\mathrm m\infty}$ rather than the schematic integer-valued functions above.
The important design choice is to isolate all infinity cases ($R-Q^j=0$) at the boundary.
For an \AEconj{} orientation those equalities are impossible, but the generic library theorem
should either assume them away or state the valuation in an extended codomain.

\subsection{Recommended proof order}

\begin{enumerate}
\item Formalize the denominator valuation
$\vp(Q^{\binom n2}\prod_{j=1}^n(Q^j-1))=m\binom n2$.
\item Formalize the finite minimum sum
$\sum_{j=1}^{n-1}\min(a,mj)$ with the threshold $q_*$.  Keep the tie correction as a
separate lemma rather than burying it in case splits.
\item Prove the off-ray factor valuations by partitioning the index range into
$j<k$, $j=k$, and $j>k$.
\item Prove the exact term ratio algebraically before any asymptotics.
\item Derive eventual strict decrease using ordinary integer inequalities; then use the
unique-minimum valuation lemma for finite sums.
\item Formalize monic-division nonexpansion and the universal Gauss bound.
\item Add the edge-operator finite spectrum and finite-strip factorization.
\item Only after the finite package is kernel-accepted, connect to Mathlib's
nonarchimedean analytic-function and power-series APIs for the identity theorem.
\end{enumerate}

This order produces useful machine-checked artifacts even if the analytic library work takes
longer.  In particular, \cref{thm:coeff-valuation,thm:universal-gauss,thm:last-term} are finite
statements with no reliance on compactness, choice of branch, or complex analysis.

\subsection{Formal status that should be claimed}

Until code is compiled by Lean with no \texttt{sorry}, the report should retain the following
status:
\begin{center}
\begin{tabular}{@{}ll@{}}
\toprule
Statement & Current status\\
\midrule
External-prime analytic rigidity & proved here, paper level\\
Exact coefficient valuation & proved here, exact Python audit\\
Universal Gauss lower bound & proved here, paper level\\
Off-ray valuation and last-term dominance & proved here, exact Python audit\\
Finite edge spectrum and strip defect & proved here, paper level\\
Robba-to-Tate and finite-rank interfaces & conditional research targets\\
Full \AEconj{} conjecture & open\\
\bottomrule
\end{tabular}
\end{center}

\section{Priority research program}
\label{sec:priorities}

\subsection{Priority 1: formalize the finite Newton wall}

The finite local package is exact, short enough for independent checking, and directly
usable in later experiments.  Formalization should include the resonant tie correction and
last-term dominance, not only the quadratic asymptotic.  This is the fastest route to a
kernel-certified contribution distinct from the supplied report.

\subsection{Priority 2: measure the edge-defect quotient}

For exact two-row interpolation polynomials, compute
\[
 S_{K,N}(z)=
\frac{P_N(\alpha z)-\lambda P_N(z)}
 {\prod_{k=K}^{N}(z-Q^k)}
\]
at external primes.  The decisive statistic is
\[
 \frac1N\log\lVert S_{K,N}\rVert_\rho
 +\log\rho
 -\frac1N\log\sigma_{p,\deg P_N}.
\]
A limit bounded away from zero confirms saturation of the finite-strip lower bound.  A limit
at or below zero would signal exactly the subexponential quotient needed by
\cref{target:defect-quotient} and should be investigated symbolically.

\subsection{Priority 3: search for finite-rank closure of defects}

The supplied report contains several exact one-axis and two-axis defect identities.  Treat
them as generators of a module under the two dilation operators, not as isolated special
functions.  Rank stabilization over $\Q(z)$ is the positive signal.  Every candidate relation
must be reconstructed exactly and then verified by symbolic polynomial identities.

\subsection{Priority 4: classify external-prime resonance}

For a candidate pair $(M,A)$, label every external prime by:
\[
 p\mid M:\quad 3\in\overline{\langle2\rangle}\subset\Z_p^\times?
 \qquad
 p\mid A:\quad 2\in\overline{\langle3\rangle}\subset\Z_p^\times?
\]
At nonresonant primes the Tate edge operator has a uniform inverse bound.  At resonant primes,
record $\mu_{p,d}$ and compare it with the degree and defect growth.  A theorem forcing at
least one nonresonant external orientation would materially strengthen the route; no such
theorem is currently known.

\subsection{Priority 5: look for an arithmetic overconvergence theorem}

The best possible bridge would start from a globally defined arithmetic object: a diagonal of
a rational function, a $G$-function, a finite-rank $q$-difference module, or another class
with controlled denominators at almost every prime.  If the exact integer samples were
certified within that class, an external prime would force positive radius and the local
rigidity theorem would finish.  The construction must be genuinely algebraic across places;
an archimedean interpolation with rational sample values is not enough.

\subsection{Explicit kill criteria}

To avoid reopening closed avenues, each program has a finite kill criterion:

\begin{longtable}{@{}p{0.30\textwidth}p{0.30\textwidth}p{0.32\textwidth}@{}}
\caption{Positive signals and kill criteria.}\label{tab:kill-criteria}\\
\toprule
Program & Positive signal & Kill criterion\\
\midrule
\endfirsthead
\toprule
Program & Positive signal & Kill criterion\\
\midrule
\endhead
Defect quotient & $\log\lVert S_{K,N}\rVert_\rho=o(N)$ after the finite spectral loss & exact growth $\ge(c+o(1))N$ with $c> -\log\rho$ in every admissible construction\\
Finite-rank cocycle & rank over $\Q(z)$ stabilizes & exact rank grows unboundedly with truncation order\\
External resonance & one external orientation has $\sigma_p>0$ & all external orientations are resonant, with small divisors matching every gain\\
Overconvergence & one number-field series has positive radius and place-independent samples & coefficient denominators have quadratic escape at every external orientation\\
Product formula & reduced residual has matching quadratic archimedean decay & reduced height remains quadratic while decay is only linear\\
\bottomrule
\end{longtable}

\section{Honest bottom line}
\label{sec:bottom-line}

No proof of \cref{conj:AE} has been obtained in this report.  The Four Exponentials barrier
has not been bypassed, and the rival team's rumored Lean development could not be verified
from public sources.

There is nevertheless a concrete advance relative to the supplied corpus.  A hypothetical
counterexample has an external nonarchimedean accumulation point, and at that point:

\begin{enumerate}
\item one geometric ray already excludes every analytic or meromorphic germ unless the two
  output integers are exact integral powers of one another;
\item one adjacent row excludes every finite-Puiseux germ by a coefficient-spectrum argument;
\item every exact finite polynomial interpolation of the native ray incurs an intrinsic
  quadratic external-prime denominator, independent of basis;
\item the canonical geometric Newton terms diverge quadratically at every off-ray unit
  multiple, and the partial sums are eventually equal in valuation to their final term;
\item every exact finite adjacent strip forces an exponentially large edge-defect quotient on
  a strict disk, up to the explicitly isolated small-divisor factor;
\item any uniformly overconvergent family of finite-prefix interpolants would compactify to the
  forbidden analytic germ.
\end{enumerate}

The remaining problem is no longer ``find some interpolation.''  Interpolation is abundant.
The sharp bridge is:

\begin{resultbox}{The new finishing interface}
Derive, from the global arithmetic relation $M=2^x$, $A=3^x$ and the integer semigroup alone,
either
\begin{enumerate}[label=(\roman*)]
\item a place-independent positive-radius interpolation at one external prime;
\item a Robba-ring object forced to descend to a Tate or finite-Puiseux germ;
\item a finite-rank consistent two-dilation module regular at the external accumulation point;
  or
\item a subexponential finite-strip edge quotient contradicting
  \eqref{eq:defect-lower}.
\end{enumerate}
Any one of these closes the conjecture immediately.  The exact Newton-wall theorems explain
why the most obvious completion fails and quantify the denominator growth that a successful
construction must cancel.
\end{resultbox}

This is a narrower and more testable frontier than the initial statement.  It also suggests a
useful interpretation of any genuine rival breakthrough: unless it introduces new
transcendence input, it should be possible to locate its decisive step as an adelic
regularity theorem, a finite-rank difference closure, or an external-prime denominator
cancellation.  Those are now explicit interfaces rather than vague hopes.

\appendix

\section{Exact audit code}
\label{app:audit-code}

\begin{lstlisting}[style=pseudocode,language=Python]
from fractions import Fraction
from math import prod


def vp_int(n, p):
    if n == 0:
        return 10**9            # infinity sentinel; excluded in assertions
    n = abs(n)
    e = 0
    while n % p == 0:
        e += 1
        n //= p
    return e


def vp_frac(x, p):
    return vp_int(x.numerator, p) - vp_int(x.denominator, p)


def c_n(Q, R, n):
    if n == 0:
        return Fraction(1)
    num = prod(R - Q**j for j in range(n))
    den = Q**(n*(n-1)//2) * prod(Q**j - 1 for j in range(1, n+1))
    return Fraction(num, den)


def phi(Q, z, n):
    return prod(z - Q**j for j in range(n))


def predicted_c_val(Q, R, p, n):
    m = vp_int(Q, p)
    a = vp_int(R, p)
    if n == 0:
        return 0
    if a == 0:
        return vp_int(R - 1, p) - m*n*(n-1)//2
    qstar = (a - 1)//m
    epsilon = 0
    if a % m == 0:
        h = a//m
        if n - 1 >= h:
            epsilon = vp_int(R - Q**h, p) - a
    if n < qstar + 2:
        return (sum(vp_int(R - Q**j, p) for j in range(n))
                - m*n*(n-1)//2)
    return (m*qstar*(qstar+1)//2
            + a*(n - 1 - qstar)
            + epsilon
            - m*n*(n-1)//2)


def predicted_phi_val(Q, u, k, p, n):
    m = vp_int(Q, p)
    if n <= k:
        return vp_int(phi(Q, u*Q**k, n), p)
    return m*(k*n - k*(k+1)//2) + vp_int(u - 1, p)


cases = [
    (35, 22, 5, 2, 2, "unit-output case"),
    (10, 110, 5, 2, 1, "resonant valuation tie"),
    (25, 75, 5, 2, 1, "integral slope, unit mismatch"),
]

for Q, R, p, u, k, label in cases:
    partial = Fraction(0)
    for n in range(13):
        coeff = c_n(Q, R, n)
        basis = phi(Q, u*Q**k, n)
        term = coeff*basis
        partial += term
        assert vp_frac(coeff, p) == predicted_c_val(Q, R, p, n)
        assert vp_int(basis, p) == predicted_phi_val(Q, u, k, p, n)
        print(label, n, vp_frac(coeff, p), vp_int(basis, p),
              vp_frac(term, p), vp_frac(partial, p))
\end{lstlisting}

\section{Proof details for monic division and weighted Gauss norms}
\label{app:gauss-details}

For completeness, this appendix records the two standard norm facts used in
\cref{thm:universal-gauss,thm:finite-strip}.

\begin{lemma}[Monic division is nonexpanding]\label{lem:monic-division}
Let $K$ be nonarchimedean, let $G\in\mathcal O_K[z]$ be monic, and divide
$P=QG+R$ with $\deg R<\deg G$.  Then
\[
 \lVert Q\rVert_{\mathrm G},\lVert R\rVert_{\mathrm G}
 \le\lVert P\rVert_{\mathrm G}.
\]
\end{lemma}

\begin{proof}
Run long division from the highest degree downward.  At each step the quotient coefficient is
the current leading coefficient, whose norm is at most the current polynomial norm.  Since
all coefficients of $G$ have norm at most one, subtracting that coefficient times a shift of
$G$ cannot increase the norm by the ultrametric inequality.  Induction proves the claim.
\end{proof}

\begin{lemma}[Norm of a linear factor]\label{lem:linear-factor-norm}
For the weighted Gauss norm \eqref{eq:weighted-gauss},
\[
 \lVert z-a\rVert_\rho=\max(\rho,|a|).
\]
The weighted Gauss norm is multiplicative on polynomials, so
\[
 \left\lVert\prod_i(z-a_i)\right\rVert_\rho
 =\prod_i\max(\rho,|a_i|).
\]
\end{lemma}

\begin{proof}
The first formula is immediate from the two coefficients.  Multiplicativity is the Gauss
lemma for a nonarchimedean norm, after rescaling the variable when $\rho$ lies in the value
group.  For arbitrary real $\rho>0$, extend scalars to a complete algebraic closure, whose
value group is dense in $\mathbb R_{>0}$, approximate $\rho$ by value-group radii, and use
continuity of the finite maximum defining the polynomial norm.  Only the polynomial case is
used here.
\end{proof}

\section{Compact theorem inventory}
\label{app:inventory}

\begingroup\small
\begin{longtable}{@{}p{0.32\textwidth}p{0.18\textwidth}p{0.42\textwidth}@{}}
\caption{Inventory of the new package.}\label{tab:inventory}\\
\toprule
Statement & Status & Mathematical role\\
\midrule
\endfirsthead
\toprule
Statement & Status & Mathematical role\\
\midrule
\endhead
Geometric-ray rigidity, \cref{thm:ray-rigidity} & proved here & Classifies every analytic germ on a contracting geometric ray.\\
Accumulating edge law, \cref{thm:edge-law} & proved here & Converts adjacent-row data into an exact Taylor-spectrum condition.\\
Meromorphic/Puiseux exclusion, \cref{thm:puiseux} & proved here & Shows finite ramification cannot hide the nonintegral exponent.\\
Exact coefficient law, \cref{thm:coeff-valuation} & proved and audited & Gives the full quadratic local denominator, including ties.\\
Universal Gauss bound, \cref{thm:universal-gauss} & proved here & Makes the denominator loss independent of interpolation basis and degree.\\
Off-ray valuation, \cref{lem:phi-offray} & proved and audited & Computes the exact local cost of moving from the native ray.\\
Term ratio and last-term dominance, \cref{prop:term-ratio,thm:last-term} & proved and audited & Establishes maximal $p$-adic divergence and adjacent-row repulsion.\\
Finite edge spectrum, \cref{prop:finite-spectral} & proved here & Isolates the exact small-divisor loss.\\
Finite-strip defect, \cref{thm:finite-strip} & proved here & Forces a large quotient behind every long exact adjacent strip.\\
No overconvergent family, \cref{thm:no-overconvergent-family} & proved here & Turns any uniform radius gain into the forbidden analytic germ.\\
Robba-to-Tate descent, \cref{target:robba-tate} & open target & Most direct local finishing bridge.\\
Finite-rank edge cocycle, \cref{target:finite-rank-cocycle} & open target & Connects the arithmetic grid to independent-$q$ rigidity theory.\\
Subexponential defect quotient, \cref{target:defect-quotient} & open target & Finite, computationally testable finishing criterion.\\
\bottomrule
\end{longtable}
\endgroup

\begin{thebibliography}{99}
\small

\bibitem{AE1944}
L.~Alaoglu and P.~Erd\H{o}s,
\newblock \emph{On highly composite and similar numbers},
\newblock Transactions of the American Mathematical Society \textbf{56} (1944), 448--469.
\newblock \href{https://doi.org/10.1090/S0002-9947-1944-0011087-2}{doi:10.1090/S0002-9947-1944-0011087-2}.

\bibitem{Unified2026}
V.~Reshetnikov and collaborators,
\newblock \emph{The Alaoglu--Erd\H{o}s Integer-Exponent Conjecture: Machine-Checked Partial Results, Exact Conditional Structure, Determinant Methods, and the Remaining Barrier},
\newblock working report supplied as \texttt{alaoglu-erdos-unified-report.tex}, 22 August 2026.

\bibitem{Waldschmidt2023}
M.~Waldschmidt,
\newblock \emph{The Four Exponentials Problem and Schanuel's Conjecture},
\newblock in \emph{Mathematics Going Forward}, Lecture Notes in Mathematics, Springer, 2023,
579--592.
\newblock \href{https://doi.org/10.1007/978-3-031-12244-6_39}{doi:10.1007/978-3-031-12244-6\_39}.

\bibitem{Robert2000}
A.~M. Robert,
\newblock \emph{A Course in $p$-adic Analysis},
\newblock Graduate Texts in Mathematics 198, Springer, 2000.

\bibitem{Schikhof1984}
W.~H. Schikhof,
\newblock \emph{Ultrametric Calculus: An Introduction to $p$-Adic Analysis},
\newblock Cambridge Studies in Advanced Mathematics 4, Cambridge University Press, 1984.

\bibitem{BGR1984}
S.~Bosch, U.~G\"untzer, and R.~Remmert,
\newblock \emph{Non-Archimedean Analysis: A Systematic Approach to Rigid Analytic Geometry},
\newblock Grundlehren der mathematischen Wissenschaften 261, Springer, 1984.

\bibitem{GasperRahman}
G.~Gasper and M.~Rahman,
\newblock \emph{Basic Hypergeometric Series}, second edition,
\newblock Encyclopedia of Mathematics and its Applications 96, Cambridge University Press, 2004.

\bibitem{Bezivin1992}
J.-P. B\'ezivin,
\newblock \emph{Sur les \`equations fonctionnelles aux $q$-diff\'erences},
\newblock Aequationes Mathematicae \textbf{43} (1992), 159--176.

\bibitem{BezivinBoutabaa1992}
J.-P. B\'ezivin and A.~Boutabaa,
\newblock \emph{Sur les \`equations fonctionnelles $p$-adiques aux $q$-diff\'erences},
\newblock Collectanea Mathematica \textbf{43} (1992), 125--140.

\bibitem{SchaefkeSinger2019}
R.~Sch\"afke and M.~F. Singer,
\newblock \emph{Consistent systems of linear differential and difference equations},
\newblock Journal of the European Mathematical Society \textbf{21} (2019), 2751--2792.
\newblock Preprint: \url{https://arxiv.org/abs/1605.02616}.

\bibitem{AnthropicFable}
Anthropic,
\newblock \emph{Claude Fable 5 and Claude Mythos 5}, 9 June 2026,
\newblock \url{https://www.anthropic.com/news/claude-fable-5-mythos-5}, accessed 23 August 2026.

\bibitem{TaoJacobian}
T.~Tao,
\newblock \emph{A digestion of the Jacobian conjecture counterexample}, 21 July 2026,
\newblock \url{https://terrytao.wordpress.com/2026/07/21/a-digestion-of-the-jacobian-conjecture-counterexample/}, accessed 23 August 2026.

\bibitem{ICARM30}
ICARM Elliptic Curve Rank Leaderboard,
\newblock \emph{Curve 273}, August 2026,
\newblock \url{https://elliptic-rank.icarm.cloud/curve/273}, accessed 23 August 2026.

\bibitem{MathWorldRank30}
E.~W. Weisstein,
\newblock \emph{Elliptic Curve Rank}, MathWorld, updated August 2026,
\newblock \url{https://mathworld.wolfram.com/EllipticCurveRank.html}, accessed 23 August 2026.

\bibitem{EpochHadamard}
Epoch AI,
\newblock \emph{Hadamard Matrix of Order 668}, FrontierMath: Open Problems, updated 12 August 2026,
\newblock \url{https://epoch.ai/frontiermath/open-problems/hadamard}, accessed 23 August 2026.

\bibitem{OEISHadamard}
OEIS Foundation,
\newblock \emph{A007299: orders associated with the Hadamard matrix existence record}, updated 2026,
\newblock \url{https://oeis.org/A007299}, accessed 23 August 2026.

\end{thebibliography}

\end{document}
